\documentclass[../main.tex]{subfiles}

\begin{document}


\vfill
\newpage
\chapter{Change of Coordinates}
\section{Chain Rule in Tensor Notation}

The core is: \textit{\textbf{contract}} the notation!

 Consider a mapping from independent variables $\mu, \nu$ to one dependent variable to another.

Using chain rule, 
$$\Pfrac{f}{\mu}=\Pfrac{F}{a}\Pfrac{A}{\mu}+\Pfrac{F}{b}\Pfrac{B}{\mu}+\Pfrac{F}{c}\Pfrac{C}{\mu}$$

and

$$
\Pfrac{f}{\nu}=\Pfrac{F}{a}\Pfrac{A}{\nu}+\Pfrac{F}{b}\Pfrac{B}{\nu}+\Pfrac{F}{c}\Pfrac{C}{\nu}
$$

From here, we use \textit{Einstein Summation Notation} to write them simply. 
$$\Pfrac{F}{a^i}\Pfrac{A^i}{\mu}$$
Finally, we also contract $\mu$ and $\nu$ to $\mu^i$ s.t.,

\begin{equation}
\Pfrac{f}{\mu^{\alpha}}=\Pfrac{F}{a^i}\Pfrac{A^i}{\mu^{\alpha}}
\end{equation}

\section{Inverse Functions of Several Variables}

Consider the mapping from $\Z$ to $Z'$ and then again to $Z$. By the chain rule, we have:

\begin{equation}
\Pfrac{Z^i(Z'(Z))}{Z^j}=\Pfrac{Z^i}{Z^{i'}}\Pfrac{Z^{i'}}{Z^j}.=J^i_{i'}J^{i'}_j
\end{equation}

Since differention of $Z^i$ with respect to $Z^j$ is $\delta^i_j$, we have:

\begin{equation}
    J^i_{i'}J^{i'}_j=\delta^i_j
\end{equation}

(1.3) means that $J^i_{i'}$ and $J^{i'}_j$ are inverses of each other. s.t.,
$$(J^\cdot_{\cdot'})^{-1}=J^{\cdot'}_\cdot.$$

\end{document}